\section{Vocabulary}

The report uses a small vocabulary throughout. Each term is defined here, in the
order a newcomer meets them, so that no later section has to introduce it.

\subsection{What is compared}

\begin{description}[leftmargin=3.1cm, style=nextline]
\item[Policy (politique)] One way of carrying out an audit. The programme
compares two: a declared automation and an agent driven by a language model.

\item[Arm (bras)] One concrete implementation of a task inside a family. The
arms share the briefing and the journal; only the interpreter changes. A family
is typically run with four: the declared rules, a script reading by convention,
the model, and a witness.

\item[D1] The declared automation: the rules written explicitly for the task.
\item[D1+] D1 \emph{augmented with every rule derivable from the public inputs}.
This is the reference a model must beat for a result to mean anything. A
competent reader who knows the briefing must not be able to state the rule; if
they can, D1+ would have it, and the family measures nothing.

\item[Family (famille)] A group of variants testing the same competence. Its
variants count as a single observation, because they test one thing. Two
families share no simulator, no vocabulary and no failure mechanism.

\item[Closure test (test de clôture)] The check performed before measuring a
family: a competent automation must fail. If it succeeds, the family measures
nothing about the model and is rejected. The test must be executable code, not a
doctrine.

\item[Filter (filtre préalable)] The cheaper check that precedes construction: if
the rule solving the case can be written in one line, a script will have it.
\end{description}

\subsection{What makes a result verifiable}

The chain has four links, and each section of this report describes one of them.

\begin{description}[leftmargin=3.1cm, style=nextline]
\item[Briefing (public)] Everything given to an arm before it runs: the task, the
target, the tools. It never contains the posture or the expected verdict.

\item[Oracle (private)] The expected answer, written before the run and never
shown to the arms. It states which state was installed and which fact is
observable, never which decision the policy should return.

\item[Witness (témoin)] A fresh value drawn when the fixture starts, whose value
appears nowhere in the briefing. It proves that an observed disclosure is fresh
and not a stale artefact.

\item[Journal] The archived log of every tool call made during a run, with its
arguments, its result and its reference. It is what the evaluator recounts; the
policy's own conclusion is not trusted.

\item[Evaluator (offline)] The program that rereads the journal and recomputes
the outcome. A cited piece of evidence absent from the log is rejected.

\item[Fixture] The state installed on the target before a trial: a healthy
posture, a vulnerable one, or a frozen corpus read offline.
\end{description}

\subsection{How results are judged}

\begin{description}[leftmargin=3.1cm, style=nextline]
\item[Funnel] The three stages of a finding, in order: \texttt{candidates},
\texttt{filtered}, \texttt{confirmed}. Only the last counts as a true positive.

\item[Strict-v3] The normalised vulnerability contract and the metrics derived
from it. The \texttt{confirmed.*} metrics count only findings whose evidence has
been independently validated against the tool log.

\item[Campaign admission] The set of conditions a campaign must meet before its
results can be aggregated. When one is unsatisfied, the results remain
descriptive and are declared as such.

\item[V0] The validity gate placed before every hypothesis: a control task that
is solvable without inference. If either arm fails it, no downstream result is
readable.

\item[H1 to H6] The six hypotheses of the register, each naming a kind of
competence a model might have where a script does not. H2 is the only one
measured; H4 is the one the rejected family attempted.

\item[C2] A candidate scenario rejected after two adversarial reviews, on paper,
before any run. Its rejection is a published negative result.

\item[Posture] The state of a fixture with respect to security: healthy (no flaw
installed) or vulnerable (flaw installed).
\end{description}
